
\vspace{-1mm}
\section{Discussion}

\subsection{Limitations}
\vspace{-1mm}
Our datasets and prompts are simple, and it is unclear how important this effect would be for more realistic datasets and more agentic settings.

In some of our settings, LLMs may infer that they are not rating actions that they generated: while most of our settings are using generations that are generated by the LLMs that we use during rating, this is not the case for the computer-use setting. Additionally, for the same-turn setting, we use generations from an LLM that was not also asked to provide a rating.


% fabien feedback:
% \subsection{Implications}

% [Walking thorugh why this is an issue when setting threshsold.]

% [Walking thourgh why this is a good interp challegne / is maybe deception related / sth like that. And how this is an issue for collusion between monitor and policy --> collusion might be an average-case failure, not sth that requires scheming?]
\subsection{Implications}

\paragraph{Implications for building agent monitors.} %Self-attribution bias has immediate implications for the design of agentic systems that rely on model-based self-monitoring.
%Many deployed pipelines use a fixed threshold on a model’s own evaluation for approving a code change if a self-assessed correctness score exceeds a cutoff, or executing a tool action if an estimated risk score falls below a safety threshold. Our results show that such thresholds are systematically unreliable when evaluation is on-policy: self-attribution selectively inflates scores for incorrect or unsafe actions, precisely in the regime where accurate discrimination is most important. 
While it might be tempting to use prompt format like follow-up questions because of their simplicitly of implementations or because of the cost savings that sharing a KV cache with the main agent can offer, using this sort of prompt format likely comes at the cost of reduced performance. Additionally, noticing this reduced performance might be difficult since this reduced performance is often only visible on-policy. Developers using this kind of prompt format that use this prompt format despite its performance drawbacks should avoid only relying on static off-policy evaluations.

\paragraph{Implications for AI control and alignment research.} Our findings suggest that self-monitoring failures that have similar properties to collusion failures can arise even on models that are very unlikely to be explicitly scheming. Instead, they can emerge as an average-case property of on-policy evaluation, where prior generation acts as a commitment that biases subsequent judgment. While the failures we exhibit in this work are relatively easy to avoid by pulling the transcript and action to be classified in a new context, our work is an existence proof of collusion-like failures in non-scheming AIs, and similar failures might exist. Our work shows that ruling out scheming is not sufficient to rule out collusion-like failures, and that finding collusion-like failures may not be decisive evidence of scheming.

Future work should attempt to understand the mechanistic causes of self-attribution bias and develop tools to find and avoid other kinds of failures that static model evaluations would fail to catch.
%  In this sense, self-attribution bias represents a structural form of monitor–policy collusion: the same model both proposes an action and then evaluates it, and the evaluation becomes systematically correlated with the proposal in a self-favoring direction. Importantly, this failure mode does not require the model to represent or pursue misaligned goals; it arises even in benign settings and for otherwise well-calibrated evaluators.
% This distinction matters for interpretability and safety research. Many existing approaches to detecting or mitigating collusion focus on worst-case behaviors such as explicit deception or hidden reasoning. Our results indicate that even absent such behaviors, on-policy self-monitoring can fail in a predictable and systematic way. This suggests that safeguards based solely on training models to “be honest” or to avoid deception may be insufficient. Instead, architectural separation between generation and evaluation—such as off-policy auditing, cross-model monitoring, or delayed evaluation on frozen artifacts—may be necessary to preserve reliable oversight.

% Our results raise questions about the use of language models as general-purpose judges in autonomous systems. While prior work shows that models can be useful evaluators of externally supplied artifacts, our findings highlight a sharp dissociation between off-policy evaluation and on-policy self-monitoring. This suggests that evaluation performance measured on static benchmarks may substantially overestimate reliability in agentic settings, where models are asked to assess their own recent actions. Designing evaluation protocols that respect this distinction is therefore crucial for the safe deployment of autonomous language model agents.




\vspace{-1mm}
\section{Conclusion}
\vspace{-1mm}
We show that LLMs display \textit{self-attribution bias}: they rate their own actions as safer and more correct than identical content under neutral framing. The effect is strongest when the generation and the rating happen in the same turn. Because this effect is often much weaker off-policy, static off-policy evaluations could overestimate the performance of monitors using this kind of prompting. %The effect is strongest in high stakes same-turn settings, where models systematically downplay the risks of actions they have just taken. %This post-hoc rationalization is distinct from static preferences—it emerges precisely when stable judgment is most critical.
%These findings challenge the assumption that LLMs can serve as reliable self-evaluators in domains ranging from scalable oversight to autonomous agents. If models inflate their own safety and correctness assessments, self-refinement pipelines risk amplifying errors rather than correcting them. Future work should explore mitigation strategies, including external evaluators, attribution-aware prompting, and training interventions that disentangle self-assessment from self-justification.
Our results highlight a tension in model design: as LLMs are trained to behave like coherent agents, their coherence can distort evaluations. Developers should be careful with prompt formats and off-policy evaluations that let LLMs infer they are evaluating themselves. 
%Specifically,
%we encourage the community to only use off-policy evaluations, or use conservative thresholds in self-monitoring deployments.

%Addressing this attribution bias is therefore essential for building trustworthy AI systems.
